\documentclass[10pt]{article}
\usepackage[T1]{fontenc}
\usepackage[utf8]{inputenc}
\usepackage[margin=0.7in]{geometry}
\usepackage{graphicx}
\usepackage{booktabs}
\usepackage{amsmath}
\usepackage{amssymb}
\usepackage{array}
\usepackage{caption}
\captionsetup{font=small}
\usepackage[hidelinks]{hyperref}
\usepackage{titlesec}
\titleformat*{\section}{\large\bfseries}
\titleformat*{\subsection}{\normalsize\bfseries}
\setlength{\parskip}{0.15em}
\graphicspath{{./}{figures/}}

\title{\textbf{Supplementary Information: An Autonomous Medicinal-Chemistry Scientist for Payload-Aware Linker Design}}
\author{}
\date{}
\begin{document}
\maketitle
\section{Full reasoning chains}
\subsection{Cytotoxin}
Query: \emph{cytotoxic payload ADC linker cleavable Val-Cit-PABC cathepsin B protease plasma stability lysosomal release maleimide}. Retrieved 10 passages; 16/16 exemplars grounded across 1 source(s); confidence 0.76 (high).

Derived rule: cleavage=\textbf{protease-cleavable}, rigidity=semi-rigid, stability=high. Rationale: Among the 16 exemplars, protease-cleavable linkers (mc-ValCitPABC, mc-cBuCitPABC, GluValCitPABC, AspValCitPABC, mc-GGFG) dominate with 8 instances across diverse cytotoxic payloads (MMAE, PBD dimer, Dxd, PE38), and these are the only linker class explicitly associated with high or very high plasma stability. Non-cleavable linkers (mc, MCC) are used selectively with payloads that retain activity as amino-acid conjugates (MMAF, DM1), while acid-cleavable linkers (hydrazone, carbonate) appear for specific payloads (calicheamicin, SN-38) but lack reported stability data. The addition of charged re
\begin{table}[h]\centering\scriptsize\begin{tabular}{l l l l}\toprule Payload & Linker & Cleavage & Source \\ \midrule MMAE & mc-ValCitPABC & protease-cleavable & biomedicines-11-03080 \\
Calicheamicin & 4-(4-acetylphenoxy) bu & acid-cleavable & biomedicines-11-03080 \\
Maytansine DM1 & MCC & non-cleavable & biomedicines-11-03080 \\
Calicheamicin & Hydrazone & acid-cleavable & biomedicines-11-03080 \\
Deruxtecan (Dxd) & mc-GGFG-aminomethoxy & protease-cleavable & biomedicines-11-03080 \\
SN-38 & mc-PEG-carbonate & acid-cleavable & biomedicines-11-03080 \\ \bottomrule\end{tabular}\end{table}
\subsection{Oligonucleotide}
Query: \emph{antibody siRNA oligonucleotide conjugate ARC linker non-cleavable rigid sulfo-SMCC stability gene silencing}. Retrieved 10 passages; 3/3 exemplars grounded across 1 source(s); confidence 0.80 (high).

Derived rule: cleavage=\textbf{non-cleavable-rigid}, rigidity=rigid, stability=high. Rationale: All three exemplars use non-cleavable thiol-maleimide linkers, but the rigid sulfo-SMCC linker outperformed both the flexible EMCS and the short AMAS linkers in producing the desired DAR2 species, indicating that rigidity is critical for efficient conjugation of bulky siRNA payloads. The sulfo group on SMCC likely aids solubility during conjugation of the hydrophilic oligonucleotide payload. However, plasma stability remains a concern (t1/2 \textasciitilde{}8 h in mice even with the best linker), suggesting that further optimization beyond these three linkers may be needed.
\begin{table}[h]\centering\scriptsize\begin{tabular}{l l l l}\toprule Payload & Linker & Cleavage & Source \\ \midrule siRNA & sulfo-SMCC & non-cleavable & sirna \\
siRNA & EMCS & non-cleavable & sirna \\
siRNA & AMAS & non-cleavable & sirna \\ \bottomrule\end{tabular}\end{table}
\subsection{Immunomodulator}
Query: \emph{immune stimulating antibody conjugate ISAC TLR7 TLR8 agonist imidazoquinoline linker cleavable plasma stability cytokine}. Retrieved 10 passages; 6/6 exemplars grounded across 1 source(s); confidence 0.64 (medium).

Derived rule: cleavage=\textbf{protease-cleavable}, rigidity=semi-rigid, stability=standard. Rationale: For TLR7 agonist immunomodulator payloads, the evidence shows that protease-cleavable linkers (mcValCitPABC) outperform non-cleavable linkers (mc), as anti-HER2\_(DG/Cys)\_mcE104 with a non-cleavable linker exhibited significantly lower potency than its cleavable counterpart. Multiple clinical-stage ADCs (BDC-1001, NJH-395) used non-cleavable linkers but showed limited clinical efficacy, further supporting the cleavable approach. Site-specific Cys conjugation on deglycosylated antibodies is preferred over transglutaminase-based Q295 conjugation, as the latter (DAR2 DG/Q295 mcValCitPABC-E104) sho
\begin{table}[h]\centering\scriptsize\begin{tabular}{l l l l}\toprule Payload & Linker & Cleavage & Source \\ \midrule E104 (TLR7 agonist & mcValCitPABC & protease-cleavable & isac \\
E104 (TLR7 agonist & mc (noncleavable) & non-cleavable & isac \\
TLR7 agonist & noncleavable linker & non-cleavable & isac \\
E104 (TLR7 agonist & mcValCitPABC & protease-cleavable & isac \\
TLR7 agonist & noncleavable linker & non-cleavable & isac \\
TLR7 agonist & noncleavable linker & non-cleavable & isac \\ \bottomrule\end{tabular}\end{table}
\section{Leave-one-paper-out details}
\subsection{Cytotoxin}
Withheld: biomedicines-11-03080 (leaked: False). Sources seen: pharmacology, songli2021. Predicted cleavage=\textbf{protease-cleavable}, rigidity=semi-rigid, confidence 0.841. Withheld paper states: \emph{Lysosomal cathepsin-B-cleavable Val-Cit/Val-Ala peptide linkers enable traceless intracellular release of the cytotoxin (Balamkundu \& Liu, Biomedicines 2023).} Verdict: \textbf{AGREE}.

Agent rationale (blind): The exemplars strongly favor protease-cleavable linkers for cytotoxin payloads, but with critical caveats around plasma stability: standard dipeptide linkers (Val-Ala, Val-Cit) hydrolyze within 1 hour, whereas engineered variants (sulfatase-cleavable, pyrophosphate diester, 7-AHC-based dipeptide) achieve >7 days stability. This indicates that for cytotoxins, cleavable linkers are preferred over non-cleavable ones (Kadcyla's non-cleavable linker showed inferior efficacy), but the linker must be s
\subsection{Immunomodulator}
Withheld: Immuno (leaked: False). Sources seen: isac, songli2021. Predicted cleavage=\textbf{protease-cleavable}, rigidity=semi-rigid, confidence 0.674. Withheld paper states: \emph{Imidazoquinoline TLR7/8 ISAC linkers may be cleavable to release the agonist in the tumour, but demand maximal plasma stability to avoid systemic cytokine toxicity (imidazo[4,5-c]quinoline ISAC paper).} Verdict: \textbf{AGREE}.

Agent rationale (blind): The exemplars show that protease-cleavable mcValCitPABC linkers with cysteine conjugation yielded superior potency compared to noncleavable mc linkers for TLR7 agonist payloads (E104), where the noncleavable variant exhibited 'significantly lower potency.' This suggests that efficient intracellular release of the free immunomodulatory payload is critical for activity, favoring cathepsin-cleavable dipeptide (ValCit) linkers. While clinical-stage ISACs (BDC-1001, NJH-395) used noncleavable linkers
\subsection{Oligonucleotide}
Withheld: siRNA (leaked: False). Sources seen: biomedicines-11-03080, pharmacology, songli2021. Predicted cleavage=\textbf{non-cleavable-flexible}, rigidity=flexible, confidence 0.007. Withheld paper states: \emph{Antibody-siRNA conjugates favour a rigid, NON-cleavable linker (sulfo-SMCC); a cleavable linker that prematurely sheds the polyanionic oligonucleotide is a liability (ARC paper, Bioconjugate Chem. 2025).} Verdict: \textbf{AGREE}.

Agent rationale (blind): No exemplars were provided for the oligonucleotide payload class, so no evidence-based rule can be derived. The recommendations above are speculative defaults based on general ADC and oligonucleotide conjugation principles: oligonucleotides are large, hydrophilic, negatively charged payloads that typically benefit from flexible, non-cleavable linkers with PEG spacers to maintain solubility and accommodate their bulk. Without concrete literature exemplars, confidence is extremely low.
\section{Candidate dossiers}
\subsection{Cytotoxin}
\noindent\textbf{Cytotoxin \#1} \\ \texttt{Nc1c(C(=O)SSc2ccccn2)nnn1CC(=O)NCc1ccccc1} \\ Rule fit: This linker employs a disulfide trigger (pyridyldisulfide) intended for intracellular glutathione-mediated cleavage, which partially aligns with the cleavable preference for MMAE payloads but contradicts the derived rule favoring protease-cleavable (ValCit-type) mechanisms over redox-cleavable ones. The triazole core and benzylamide sidec \\ Nearest clinical: SPDB disulfide (Tc 0.339, mirvetuximab soravtansine (Elahere)). Steps 6, est.\ cost \$2400, 15.0 d, P(success) 0.6. \\ Route: Synthesis of 5-amino-1H-1,2,3-triazole-4-carboxylic acid via Dimroth triazole formation from appropriate azide and cyanoacetamide precursors; N-alkylation of triazole nitrogen with bromoacetyl bromide to install the chloroacetamide handle; Amide coupling of bromoacetyl-triazole intermediate with benzylamine (HATU/DIPEA \\ Failure modes: Premature disulfide reduction in plasma by albumin thiols and circulating glutathione, leading to early payload release and off-target toxicity; Poor DAR homogeneity due to competing disulfide scrambling during conjugation to antibody cyste. Validation: Plasma stability assay in human, mouse, and cynomolgus plasma (LC-MS monitoring of intact linker-payload vs free MMAE, t1/2 determination at 37\textdegree{}C over 96 h); Glutathione-mediated cleavage kinetics (10 mM GSH, pH 7.4 and pH 5.5, LC-MS quanti. \\[0.4em]
\noindent\textbf{Cytotoxin \#2} \\ \texttt{O=C(Cn1cc(C(=O)NCC(=O)SSc2ccccn2)nn1)NCc1ccccc1} \\ Rule fit: This linker employs a disulfide trigger (pyridyldisulfide) intended for intracellular glutathione-mediated cleavage, which partially aligns with the cleavable preference for MMAE but contradicts the derived rule favoring protease-cleavable (ValCit-type) linkers that offer superior plasma stability. The triazole and benzylamine moieties pr \\ Nearest clinical: SPDB disulfide (Tc 0.344, mirvetuximab soravtansine (Elahere)). Steps 7, est.\ cost \$2700, 17.5 d, P(success) 0.6. \\ Route: CuAAC or thermal cycloaddition to form the 1,2,3-triazole ring from azidoacetic acid derivative and propargyl amide precursor; Amide coupling of triazole-acetic acid with benzylamine (HATU/DIPEA); Amide coupling of triazole-3-carboxylic acid with glycine derivative (HATU/DIPEA); Activation of glycine carboxyl as thiol  \\ Failure modes: Premature disulfide reduction by plasma thiols (albumin Cys34, glutathione) causing payload release in circulation; Competing disulfide scrambling with antibody interchain disulfides leading to heterogeneous DAR and aggregation; Insufficien. Validation: Plasma stability assay in human, mouse, and cynomolgus plasma (LC-MS/MS monitoring of free MMAE release, t1/2 target >72 h); Glutathione-mediated cleavage kinetics (10 mM GSH, 37\textdegree{}C, LC-MS time-course) to confirm intracellular trigger mechan. \\[0.4em]
\noindent\textbf{Cytotoxin \#3} \\ \texttt{O=C(CNC(=O)c1nc[nH]c1C(=O)SSc1ccccn1)NCc1ccccc1} \\ Rule fit: This linker employs a disulfide trigger (pyridyl disulfide--imidazole motif) intended for intracellular glutathione-mediated release of MMAE, partially aligning with the cleavage preference for conditional release but diverging from the preferred protease-cleavable mechanism (ValCit/PABC). The bis-amide backbone with aromatic imidazole and \\ Nearest clinical: SPDB disulfide (Tc 0.339, mirvetuximab soravtansine (Elahere)). Steps 6, est.\ cost \$2400, 15.0 d, P(success) 0.6. \\ Route: Synthesis of 4,5-imidazoledicarboxylic acid mono-amide via selective mono-activation and amide coupling with benzylamine (HATU/DIPEA); Amide coupling of remaining imidazole carboxylic acid with glycinamide spacer (EDC/HOBt); Activation of imidazole thiol or installation of thiol via reduction of a protected precursor;  \\ Failure modes: Premature disulfide reduction in plasma by serum albumin thiols and circulating glutathione, leading to off-target MMAE release and systemic toxicity; Competing disulfide scrambling during conjugation to antibody interchain cysteines, yield. Validation: Plasma stability assay in human, mouse, and cynomolgus plasma (LC-MS/MS quantification of free MMAE release, t1/2 target >72 h); Glutathione-triggered release kinetics (10 mM GSH, 37\textdegree{}C, LC-MS monitoring of MMAE liberation over 24 h); HIC an. \\[0.4em]
\noindent\textbf{Cytotoxin \#4} \\ \texttt{O=C(NCc1ccccc1)Nc1n[nH]c(NC(=O)SSc2ccccn2)n1} \\ Rule fit: This linker employs a disulfide trigger intended for intracellular glutathione-mediated release of MMAE, but the derived design rule strongly favors protease-cleavable linkers (e.g., ValCit-PABC) for MMAE payloads due to superior plasma stability and validated clinical precedent. The triazine core provides semi-rigidity consistent with th \\ Nearest clinical: SPDB disulfide (Tc 0.373, mirvetuximab soravtansine (Elahere)). Steps 7, est.\ cost \$2700, 17.5 d, P(success) 0.6. \\ Route: Synthesis of 2-pyridyl disulfide thiol via activation of 2-mercaptopyridine with Aldrithiol-2; Preparation of 3,5-diamino-1,2,4-triazole core by condensation of aminoguanidine with cyanamide; Acylation of one triazole amine with the activated disulfide-containing carboxylic acid (EDC/HOBt); Synthesis of N-benzylamine u \\ Failure modes: Premature disulfide reduction in plasma by serum albumin thiols and circulating glutathione, causing off-target MMAE release; Low Tanimoto to validated clinical linkers suggests unpredictable PK and potentially poor therapeutic index; Triaz. Validation: Plasma stability assay in human/mouse plasma measuring intact linker-payload by LC-MS/MS over 96 h (target t1/2 > 72 h); Glutathione-mediated cleavage kinetics (10 mM GSH, 37\textdegree{}C) monitored by HPLC to quantify MMAE release rate; HIC and SEC t. \\[0.4em]
\noindent\textbf{Cytotoxin \#5} \\ \texttt{Nc1nc(NCc2ccccc2)ncc1C(=O)OCC(=O)SSc1ccccn1} \\ Rule fit: This linker employs a disulfide trigger intended for intracellular glutathione-mediated cleavage, which partially aligns with the cleavable preference for MMAE payloads but contradicts the derived rule favoring protease-cleavable (ValCit-type) mechanisms that offer superior plasma stability. The aminopyrimidine-benzylamine scaffold introd \\ Nearest clinical: SPDB disulfide (Tc 0.361, mirvetuximab soravtansine (Elahere)). Steps 6, est.\ cost \$2400, 15.0 d, P(success) 0.6. \\ Route: SNAr or Buchwald coupling of 2,4-dichloropyrimidine-5-carboxylic acid with benzylamine to install the NCH2Ph substituent; Amination at C-2 of the pyrimidine with ammonia or ammonium source to give the 2-amino-4-(benzylamino)pyrimidine-5-carboxylic acid; Esterification of the pyrimidine-5-carboxylic acid with chloroacet \\ Failure modes: Premature disulfide reduction by serum albumin thiols and glutathione in plasma, causing off-target MMAE release and systemic toxicity; Ester hydrolysis of the glycolate linkage under physiological pH or by serum esterases, leading to linke. Validation: Plasma stability assay in human, mouse, and cynomolgus plasma (LC-MS/MS monitoring of intact ADC and free MMAE at 0--96 h); Glutathione-triggered release kinetics (10 mM GSH, 37\textdegree{}C, pH 7.4) with LC-MS quantification of liberated MMAE; Ester h. \\[0.4em]
\subsection{Immunomodulator}
\noindent\textbf{Immunomodulator \#1} \\ \texttt{O=C(Cn1c(=O)[nH]c(=O)n(CC(=O)N2Cc3ccccc3C\#Cc3ccccc32)c1=O)NCc1ccccc1} \\ Rule fit: This design employs a non-cleavable linker with a DBCO click-chemistry handle, which directly contradicts the derived rule favoring protease-cleavable linkers for TLR7 agonist payloads; however, it serves as a critical non-cleavable comparator to validate the cleavage\_preference hypothesis. The barbituric acid-based spacer provides semi-r \\ Nearest clinical: mc-GGFG (Tc 0.195, trastuzumab deruxtecan (Enhertu)). Steps 6, est.\ cost \$2400, 15.0 d, P(success) 0.65. \\ Route: Alkylation of barbituric acid with chloroacetyl chloride to install two N-chloroacetyl arms; Amide coupling of one chloroacetyl arm with benzylamine (EDC/HOBt or displacement); Nucleophilic displacement of second chloroacetyl arm with DBCO-amine intermediate bearing isoindoline-alkyne scaffold; Copper-free SPAAC conjug \\ Failure modes: Non-cleavable linker fails to release active R848 intracellularly, resulting in poor TLR7 activation and low potency consistent with BDC-1001/NJH-395 clinical limitations; Barbituric acid core may undergo ring-opening hydrolysis under lysos. Validation: Plasma stability assay in human/mouse plasma (LC-MS monitoring of intact linker-payload at 0, 24, 72, 168 h); In vitro TLR7 NF-$\kappa$B reporter assay comparing free R848, non-cleavable ADC, and cleavable mcValCitPABC-R848 ADC; Lysosomal degradat. \\[0.4em]
\noindent\textbf{Immunomodulator \#2} \\ \texttt{O=C(Cn1cc(C(=O)N2Cc3ccccc3C\#Cc3ccccc32)c(=O)[nH]c1=O)NCc1ccccc1} \\ Rule fit: This design employs a non-cleavable DBCO-based linker with a uracil-containing spacer, which directly contradicts the derived rule favoring protease-cleavable linkers for TLR7 agonist payloads; however, the rigid DBCO click-chemistry handle provides efficient strain-promoted azide-alkyne cycloaddition conjugation and the semi-rigid aromat \\ Nearest clinical: mc-GGFG (Tc 0.187, trastuzumab deruxtecan (Enhertu)). Steps 5, est.\ cost \$2000, 12.5 d, P(success) 0.85. \\ Route: Alkylation of uracil at N1 with bromoacetamide-benzylamine intermediate (K2CO3, DMF); Sonogashira coupling of 2-iodoaniline derivative with ethynylbenzene to construct the dibenzoazacyclooctyne (DBCO) precursor; DBCO ring closure via intramolecular lactamization under high-dilution conditions; Amide coupling of uracil  \\ Failure modes: Non-cleavable linker fails to release active R848 intracellularly, resulting in poor TLR7 engagement and low immunostimulatory potency consistent with clinical failures of non-cleavable TLR7 ADCs (BDC-1001, NJH-395); DBCO cyclooctyne ring s. Validation: SPAAC conjugation efficiency with azide-functionalized antibody assessed by HIC and SEC for DAR distribution and monomer purity; Plasma stability in human, mouse, and cynomolgus plasma over 14 days with LC-MS/MS quantification of intact lin. \\[0.4em]
\noindent\textbf{Immunomodulator \#3} \\ \texttt{NCC1CCC(Nc2ccc(S(=O)(=O)N3Cc4ccccc4C\#Cc4ccccc43)cc2S(N)(=O)=O)CC1} \\ Rule fit: This design employs a non-cleavable rigid linker with a DBCO conjugation handle, which directly contradicts the derived rule favoring protease-cleavable semi-rigid linkers for TLR7 agonist payloads. The rigid diarylsulfonamide-alkyne scaffold may restrict conformational freedom needed for efficient payload release, and the non-cleavable n \\ Nearest clinical: MCC (non-cleavable) (Tc 0.132, trastuzumab emtansine (Kadcyla)). Steps 4, est.\ cost \$1600, 10.0 d, P(success) 0.85. \\ Route: Sonogashira coupling of 2-iodoaniline derivative with ortho-ethynylbromobenzene to form the diaryl alkyne core; Sulfonylation of the aniline nitrogen with 4-nitrobenzenesulfonyl chloride followed by reduction to the corresponding amine; Sulfonamide formation on the second ring using chlorosulfonyl isocyanate or SO2NH2  \\ Failure modes: Insufficient intracellular R848 release due to non-cleavable linker requiring complete lysosomal antibody degradation, yielding amino acid-linker-R848 catabolites with reduced TLR7 agonism; Rigid diaryl alkyne scaffold may cause poor aqueou. Validation: Plasma stability assay in human/mouse plasma over 14 days with LC-MS/MS monitoring of intact ADC and free R848 levels; In vitro TLR7 NF-$\kappa$B reporter assay comparing free R848, linker-R848 catabolite (post-lysosomal digest), and intact ADC to. \\[0.4em]
\noindent\textbf{Immunomodulator \#4} \\ \texttt{NCC1CCC(NS(=O)(=O)c2ccc(S(=O)(=O)N3Cc4ccccc4C\#Cc4ccccc43)cc2)CC1} \\ Rule fit: This design employs a non-cleavable rigid linker with DBCO click-chemistry conjugation, which directly contradicts the derived rule favoring protease-cleavable, semi-rigid linkers for TLR7 agonist payloads. The bis-sulfonamide aromatic scaffold with an internal diyne provides extreme rigidity, which may limit payload release and immune ac \\ Nearest clinical: MCC (non-cleavable) (Tc 0.145, trastuzumab emtansine (Kadcyla)). Steps 4, est.\ cost \$1600, 10.0 d, P(success) 0.85. \\ Route: Sonogashira coupling of 2-iodobenzylamine derivative with 2-ethynylbenzylamine to form the diaryl alkyne core; Sulfonylation of the amine on one aryl ring with 4-(chlorosulfonyl)benzenesulfonyl chloride to install the bis-sulfonamide bridge; Sulfonamide formation on the second nitrogen with the trans-4-aminocyclohexyla \\ Failure modes: Non-cleavable linker prevents efficient intracellular R848 release, leading to poor TLR7 activation and low immunostimulatory potency; Extreme rigidity of the diaryl alkyne-bis-sulfonamide scaffold may cause aggregation or reduced conjugati. Validation: In vitro TLR7 NF-$\kappa$B reporter assay comparing free R848 vs. linker-R848 conjugate vs. intact ADC to assess payload release requirement; Plasma stability in human and mouse plasma over 14 days with LC-MS/MS monitoring of intact ADC and free R. \\[0.4em]
\noindent\textbf{Immunomodulator \#5} \\ \texttt{NCC1CCC(Nc2ccc(S(=O)(=O)N3Cc4ccccc4C\#Cc4ccccc43)cc2S(=O)(=O)O)CC1} \\ Rule fit: This design employs a non-cleavable rigid DBCO-based linker, which directly contradicts the derived rule favoring protease-cleavable (mcValCitPABC) linkers for TLR7 agonist payloads; however, its high composite score (0.78) and perfect stability (1.0) make it a valuable negative-control comparator. The rigid dibenzocyclooctyne scaffold pr \\ Nearest clinical: MCC (non-cleavable) (Tc 0.145, trastuzumab emtansine (Kadcyla)). Steps 4, est.\ cost \$1600, 10.0 d, P(success) 0.85. \\ Route: Sulfonylation of 4-aminobenzenesulfonic acid with 2-(chlorosulfonyl)isoindoline-DBCO precursor to form the diaryl sulfonamide core; Sonogashira coupling to install the internal alkyne completing the DBCO ring system; Reductive amination or Boc-deprotection to reveal the primary amine on the trans-4-aminocyclohexyl hand \\ Failure modes: Insufficient intracellular R848 release due to non-cleavable linker, leading to poor TLR7 activation consistent with clinical precedent; Sulfonic acid group may impair cell membrane permeability and lysosomal escape of the intact ADC-linker. Validation: SPAAC click conjugation efficiency to azide-functionalized antibody measured by HIC and RP-HPLC for DAR determination; Plasma stability in human and mouse plasma over 14 days with LC-MS/MS quantification of intact conjugate vs free R848; In. \\[0.4em]
\subsection{Oligonucleotide}
\noindent\textbf{Oligonucleotide \#1} \\ \texttt{O=C(NCc1ccccc1)c1cn(C(=O)N2Cc3ccccc3C\#Cc3ccccc32)c(=O)[nH]c1=O} \\ Rule fit: This linker employs a non-cleavable DBCO-based conjugation handle consistent with the non-cleavable-rigid design rule derived for siRNA payloads, where rigidity is critical for efficient conjugation of bulky oligonucleotides. The fused dibenzoazacyclooctyne (DBCO) core and alkyne-bridged biaryl system enforce a rigid scaffold that minimiz \\ Nearest clinical: mc-GGFG (Tc 0.178, trastuzumab deruxtecan (Enhertu)). Steps 5, est.\ cost \$2000, 12.5 d, P(success) 0.85. \\ Route: Sonogashira coupling of 2-iodobenzylamine derivative with 2-ethynylbenzaldehyde to form the diaryl alkyne core; Cyclization to construct the DBCO azacyclooctyne ring system under intramolecular lactamization conditions; Amide coupling of 6-aminouracil carboxylic acid with benzylamine using HATU/DIPEA; N-Acylation of ur \\ Failure modes: DBCO strain-promoted click with azide-modified siRNA may be sluggish due to steric encumbrance from the rigid uracil-benzamide substituent, leading to low DAR; The extended aromatic/alkyne system may cause aggregation or poor aqueous solubi. Validation: SPAAC conjugation efficiency with azide-functionalized siRNA monitored by SEC-HPLC and native PAGE to determine DAR distribution; Plasma stability assay (mouse and human, 37\textdegree{}C, 48 h) with LC-MS/MS quantification of intact ADC vs. free siRNA. \\[0.4em]
\noindent\textbf{Oligonucleotide \#2} \\ \texttt{Nc1nnc(C(=O)N2Cc3ccccc3C\#Cc3ccccc32)c(C(=O)NCc2ccccc2)n1} \\ Rule fit: This linker employs a non-cleavable DBCO-based conjugation handle consistent with the non-cleavable preference derived for siRNA payloads, enabling strain-promoted azide-alkyne cycloaddition (SPAAC) with azide-functionalized antibodies without requiring reducing conditions. The fused dibenzoazacyclooctyne core and flanking aromatic amide  \\ Nearest clinical: mc-GGFG (Tc 0.169, trastuzumab deruxtecan (Enhertu)). Steps 5, est.\ cost \$2000, 12.5 d, P(success) 0.85. \\ Route: Synthesis of 3-amino-1,2,4-triazine-5,6-dicarboxylic acid via condensation of aminoguanidine with oxalyl derivatives; Selective mono-amide coupling of C6-carboxylate with benzylamine (HATU/DIPEA) to install the benzylamide arm; Amide coupling of C5-carboxylate with 5,6-dihydro-11,12-didehydrodibenz[b,f]azocine (DBCO-am \\ Failure modes: DBCO cyclooctyne strain may lead to slow hydrolysis or thiol-yne side reactions in plasma, reducing effective conjugation yield; Low Tanimoto similarity (0.169) to nearest clinical analogue suggests unpredictable PK behavior and potential r. Validation: SPAAC conjugation efficiency with azide-modified trastuzumab measured by HIC and native MS to determine DAR distribution; Plasma stability assay (mouse and human, 37\textdegree{}C, 72 h) with LC-MS/MS quantification of intact ADC and free siRNA; Dynami. \\[0.4em]
\noindent\textbf{Oligonucleotide \#3} \\ \texttt{Nc1nc(N)c(C(=O)N2Cc3ccccc3C\#Cc3ccccc32)nc1C(=O)NCc1ccccc1} \\ Rule fit: This linker incorporates a rigid diaminopyrimidine core fused with a diarylacetylene (alkyne-bridged biaryl) motif, directly satisfying the 'non-cleavable-rigid' cleavage preference and high rigidity requirement derived from the sulfo-SMCC precedent for siRNA payloads. The DBCO handle enables strain-promoted azide-alkyne cycloaddition (SP \\ Nearest clinical: mc-GGFG (Tc 0.17, trastuzumab deruxtecan (Enhertu)). Steps 5, est.\ cost \$2000, 12.5 d, P(success) 0.85. \\ Route: Sonogashira coupling of 2-iodoaniline with 2-ethynylbenzaldehyde to form the diarylacetylene scaffold; Reductive amination or cyclization to construct the isoindoline (dihydroisoindole) DBCO-precursor ring system; Amide coupling of 2,4-diamino-5-pyrimidinecarboxylic acid with the isoindoline-DBCO amine (HATU/DIPEA); Am \\ Failure modes: Poor aqueous solubility due to extended aromatic/hydrophobic scaffold, complicating conjugation with hydrophilic siRNA in aqueous buffer; Steric hindrance from the rigid diaminopyrimidine-diarylacetylene core may slow SPAAC click reaction k. Validation: SPAAC conjugation efficiency with azide-functionalized siRNA monitored by SEC-HPLC and native PAGE to assess DAR distribution; Plasma stability assay (mouse and human, 37\textdegree{}C, 72 h) with LC-MS/MS quantification of intact conjugate vs. free si. \\[0.4em]
\noindent\textbf{Oligonucleotide \#4} \\ \texttt{Nc1nc(C(=O)N2Cc3ccccc3C\#Cc3ccccc32)c(N)nc1C(=O)NCc1ccccc1} \\ Rule fit: This linker employs a non-cleavable DBCO conjugation handle consistent with the non-cleavable-rigid design rule derived for siRNA payloads, where cleavable triggers are unnecessary because the intact conjugate can deliver oligonucleotide cargo upon internalization and lysosomal degradation. The fused dibenzoazacyclooctyne (DBCO) core with \\ Nearest clinical: mc-GGFG (Tc 0.172, trastuzumab deruxtecan (Enhertu)). Steps 5, est.\ cost \$2000, 12.5 d, P(success) 0.85. \\ Route: Synthesis of 2,4-diamino-5-pyrimidinecarboxylic acid via Atwal-type pyrimidine ring construction; Amide coupling of 2,4-diamino-5-pyrimidinecarboxylic acid with benzylamine (HATU/DIPEA) to install the benzyl amide arm; Synthesis of DBCO-amine core via Sonogashira coupling of 2-ethynylaniline with 2-iodobenzylamine foll \\ Failure modes: Diaminopyrimidine may undergo oxidative deamination or glucuronidation in vivo, reducing plasma stability below the target threshold; The benzyl amide arm lacks steric shielding and may be susceptible to hepatic amidase cleavage, leading to. Validation: DAR analysis by HIC and native SEC to assess conjugation efficiency and homogeneity of siRNA-ADC species; Plasma stability assay (mouse and human plasma, 37\textdegree{}C, 0-48 h) with LC-MS/MS quantification of intact conjugate vs. free siRNA; SPAAC c. \\[0.4em]
\noindent\textbf{Oligonucleotide \#5} \\ \texttt{Cn1c(C(=O)N2Cc3ccccc3C\#Cc3ccccc32)c(N)nc(C(=O)NCc2ccccc2)c1=O} \\ Rule fit: This linker incorporates a rigid dibenzoazacyclooctyne (DBCO) click-chemistry handle fused into a planar alkyne-bridged dibenzazepine core, directly satisfying the non-cleavable-rigid design rule derived for siRNA payloads. The non-cleavable amide backbone and extended aromatic scaffold minimize conformational flexibility, which the sulfo \\ Nearest clinical: mc-GGFG (Tc 0.183, trastuzumab deruxtecan (Enhertu)). Steps 5, est.\ cost \$2000, 12.5 d, P(success) 0.85. \\ Route: Sonogashira coupling of 2-iodoaniline with 2-ethynylaniline to form the diarylalkyne core; Cyclization of the diarylalkyne diamine to the dibenzoazepine scaffold under Pd-catalyzed or thermal conditions; Amide coupling of the dibenzoazepine nitrogen with 6-amino-1-methyl-2-oxo-1,2-dihydropyrimidine-4-carboxylic acid (H \\ Failure modes: Competing intramolecular cyclization or polymerization during Sonogashira/cyclization steps due to bifunctional alkyne-amine intermediates; Poor aqueous solubility of the extended aromatic scaffold limiting conjugation efficiency with hydro. Validation: SPAAC conjugation efficiency with azide-functionalized siRNA monitored by RP-HPLC and native PAGE to determine DAR distribution; Plasma stability assay (mouse and human, 37\textdegree{}C, 48 h) with LC-MS/MS quantification of intact ADC vs free siRNA; . \\[0.4em]
\section{Whole-conjugate Boltz-2 co-folds (cathepsin B)}
\begin{table}[h]\centering\small\begin{tabular}{l l c c}\toprule Conjugate & Kind & ipTM & Pred.\ $K_d$ (nM) \\ \midrule cyto-Disulf-0 & designed\_conjugate & 0.849 & 1199.14 \\
immu-DBCO-0 & designed\_conjugate & 0.927 & 43.21 \\
olig-DBCO-0 & negative\_conjugate & 0.871 & 74.09 \\
clinical-mcValCitPABC+MMAE & commercial\_conjugate & 0.746 & 44.78 \\ \bottomrule\end{tabular}\end{table}
\end{document}